% !TeX root = ../elteikthesis_hu.tex

\chapter{Bevezetés}
\label{ch:intro}

A modern full-stack fejlesztés egyik legösszetettebb feladata a valós idejű,
 többjátékos rendszerek létrehozása, ahol a kliensoldali élmény és a szerveroldali
  biztonság egyaránt kritikus. Jelen dolgozat ezen kihívásokat mutatja be
   a Snapszer kártyajáték webes adaptációján keresztül. 
   A fejezet elsőként a fejlesztés céljait rögzítő témabejelentőt tartalmazza, 
   majd részletezi a játék azon szabályait, amelyek az implementáció logikai
    alapját képezik.

\section{Témabejelentő}
A dolgozatom témája egy Online Többjátékos Snapszer Webes Játék megtervezése és implementálása. A célom, hogy egy letisztult modern játékfelületet létrehozzak, ami optimális játékélményt biztosít. A projekt bemutatja, hogy hogyan oldható meg egy valós idejű szinkronizációs játék webes környezetben.

A snapszer egy kiemelten népszerű kártyajáték itthon, melynek online verziója magas követelményeket támaszt a valós idejű interakció és a precíz játékszabály validáció terén. Míg asztali és mobil alkalmazások léteznek, webes környezetben még nincsen letisztult megvalósítása a játéknak.

A dolgozatom fő feladata, hogy gyakorlati példán keresztül bemutassa a full-stack webfejlesztés lehetőségeit. Egy olyan összetett rendszert valósít meg, amelynél elengedhetetlen a minimális késleltetéssel járó kommunikáció és a szerver oldali, biztonságos játékállapot kezelés. A választott JavaScript technológiai stack (React.js, Node.js, Express.js) azért ideális, mert egységes és hatékony alapot biztosít mind a felhasználói felület, mind a háttérben futó játékszerver komplex kihívásainak megoldásához és egyetlen nyelvet használ az egész stack.

\section{A dolgozat felépítése}

A dolgozat négy fő fejezetből, valamint az Irodalomjegyzékből, az Ábrajegyzékből és a Táblázatok jegyzékéből áll.

\vspace{1em}
\noindent\textbf{1. Bevezetés}

Itt mutatom be a témaválasztásom motivációját, a snapszer kártyajáték részletes szabályait, valamint a fejlesztés alapjául szolgáló témabejelentő is itt kapott helyet.

\vspace{1em}
\noindent\textbf{2. Felhasználói dokumentáció}

Ez egy kézikönyv a játékosok számára, amelyből megismerhetik a program futtatásához szükséges követelményeket és az alkalmazás funkcióit. Itt részletesen bemutatom a weboldal felépítését, az oldaltérképet, az automatikus meccskeresés és a játék indításának folyamatát. Továbbá a fejezet tartalmaz egy lépésről lépésre haladó útmutatót is, amely a tényleges játékmenetet és a virtuális kártyaasztal használatát mutatja be.

\vspace{1em}
\noindent\textbf{3. Fejlesztői dokumentáció}

Ez a program technikai leírása, amely az alkalmazás egy későbbi módosítása vagy bővítése esetén nyújt segítséget. Itt található a rendszer architektúrájának bemutatása, az alkalmazott webes technológiák (kliens és szerver oldal) ismertetése, az adatbázis felépítése, valamint a valós idejű többjátékos szinkronizáció megvalósításának részletei.

\vspace{1em}
\noindent\textbf{4. Összegzés}

A dolgozat lezárásaként értékelem az elvégzett munkát, összegzem a fejlesztés során szerzett tapasztalatokat, és felvázolom az alkalmazás lehetséges jövőbeli továbbfejlesztési irányait.

\vspace{1em}
\noindent\textbf{Irodalomjegyzék}

Itt hivatkozom a munkám során használt cikkekre, hivatalos technológiai dokumentációkra és eszközökre.

\vspace{1em}
\noindent\textbf{Ábrajegyzék}

A dolgozatban előforduló ábrák, diagramok és képernyőképek gyűjteménye.

\vspace{1em}
\noindent\textbf{Táblázatok jegyzéke}

A dolgozatban előforduló táblázatok gyűjteménye.

\section{A snapszer szabályai}
\label{sec:szabalyok}
A Snapszer (vagy 66) egy 20 lapos magyar kártyapaklival játszott stratégiai játék.
 A szoftveres megvalósítás során az alábbi kétjátékosos szabályrendszert vettem alapul.

\subsection{Kártyák értéke és a játék célja}
A játékban csak az Ász, a Tizes, a Király, a Felső és az Alsó lapok vesznek részt. 
A cél elsőként elérni a 66 pontot az ütések és bemondások értékéből.
\begin{compactitem}
    \item \textbf{Ász}: 11 pont, \textbf{Tizes}: 10 pont.
    \item \textbf{Király}: 4 pont, \textbf{Felső}: 3 pont, \textbf{Alsó}: 2 pont.
\end{compactitem}

\subsection{A játék menete}
A játékosok 5-5 lapot kapnak, az adu (tromf) színét a pakli tetejéről felfordított lap határozza meg.
\begin{compactitem}
    \item \textbf{Ütések}: Amíg a húzható pakli (talon) el nem fogy, nem kötelező sem színre színt tenni, sem ütni. 
    Az ütést a hívott szín nagyobb lapja vagy az adu viszi el.
    \item \textbf{Adu ereje}: Az adu szín bármely más színű lapnál erősebb, így az aduval való ütés felülbírálja a hívott színt függetlenül a lapértéktől.
    \item \textbf{Laphúzás}: Minden ütés után mindkét játékos húz egy új lapot, az ütést nyerő játékossal kezdve.
    \item \textbf{Kényszerek}: A talon elfogyása után színkényszer és ütéskényszer lép érvénybe: a hívott színnel megegyező lapot kell tenni, és ha lehetséges, felül kell ütni azt.
\end{compactitem}

\subsection{Speciális szabályok}
\begin{compactitem}
    \item \textbf{Bemondások}: Király és Felső pár esetén „húsz” (20 pont), adu színben „negyven” (40 pont) mondható be az első kijátszáskor.
    \item \textbf{Adu csere}: Az adu alsó birtokosa kicserélheti lapját a lent lévő adu lapra a laphúzás előtt, de csak amikőr a játékos van soron.
\end{compactitem}

\subsection{Végső pontozás}
A játszma végén a nyertes pontjait az ellenfél pontjai határozza meg:
\begin{compactitem}
    \item \textbf{3 pont}: Ha az ellenfélnek nincs ütése.
    \item \textbf{2 pont}: Ha az ellenfél 33 pont alatt maradt.
    \item \textbf{1 pont}: Ha az ellenfél elérte a 33 pontot.
    \item \textbf{Utolsó ütés}: Amennyiben a játék végéig egyik fél sem éri el a 66 pontot, az utolsó ütést megszerző játékos nyeri a játszmát.
\end{compactitem}

\section{Motiváció}
A projekt megvalósítását elsősorban az motiválta, hogy a snapszer egy nagy múltú és rendkívül népszerű kártyajáték, amely jelenleg nem rendelkezik modern, böngészőből azonnal játszható, minőségi adaptációval. Egy olyan platform létrehozása volt a célom, amely letöltés, telepítés és akár regisztráció nélkül, letisztult felülettel teszi elérhetővé ezt a játékot a szélesebb közönség számára, tiszteletben tartva a tradicionális szabályokat.

Szakmai oldalról egy valós idejű, többjátékos rendszer felépítése kiemelkedő kihívást jelentett, ami tökéletes lehetőséget biztosított a modern webes technológiák elmélyítésére. A kliens és szerver közötti folyamatos szinkronizáció, a biztonságos állapotkezelés és a reszponzív felület kialakítása révén a projekt komplex keretet ad a full-stack JavaScript fejlesztői képességem alkalmazásához és fejlődéséhez.